\begin{theorem}[Positive trace of nonzero positive semidefinite matrices]
    \label{thm:psd_trace_pos}
    \lean{Matrix.PosSemidef.trace_pos_of_ne_zero}
    \leanok
    A nonzero positive semidefinite matrix has positive trace: the trace is
    a positive real.
\end{theorem}

\begin{proof}
    \leanok
    The trace of a positive semidefinite matrix is a non-negative real, and
    it vanishes only if the matrix itself is zero.
\end{proof}

\begin{theorem}[Nonzero sector compressions have positive trace]
    \label{thm:mpdo_sector_compression_trace_pos}
    \lean{MPOTensor.mpo_compress_posSemidef}
    \lean{MPOTensor.mpo_compress_trace_pos}
    \lean{MPOTensor.sectorCompression_posSemidef}
    \lean{MPOTensor.sectorCompression_trace_pos}
    \leanok
    \uses{def:mpdo, thm:psd_trace_pos, def:mpdo_vertical_one_site_actions,
        def:mpdo_vertical_loops}
    Let $H^{(N)}$, $N\geq1$, be a density operator generated by an MPDO and let
    $P=P^\dagger$. Then $PH^{(N)}P\ge 0$, and if $PH^{(N)}P\ne 0$ its
    trace is a positive real. In particular this holds for the first-site
    compressions $P_1H^{(N+1)}P_1$. Applied to the sector projectors
    $P_{\alpha,k}$ of the vertical decomposition, this is the positivity
    $\tr(P_{\alpha,k}H^{(N)}P_{\alpha,k})>0$ from which the proof of
    \cite[Proposition~4.13]{Cirac2016MPDO_arXiv} concludes
    $\mu_{\alpha,k}d_\alpha>0$; the identification of the trace with
    $\mu_{\alpha,k}d_\alpha$ is Theorem~\ref{thm:mpdo_sector_trace_identity}.
\end{theorem}

\begin{proof}
    \leanok
    \uses{thm:psd_trace_pos, def:mpdo_vertical_one_site_actions}
    Compression by a Hermitian matrix preserves positive semidefiniteness,
    since $PH^{(N)}P=PH^{(N)}P^\dagger$. By
    Theorem~\ref{thm:psd_trace_pos}, a nonzero compression therefore has
    positive trace. The first-site case follows because
    $P_1=P\otimes\Id^{\otimes N}$ is Hermitian whenever $P$ is.
\end{proof}

\begin{definition}[Single-site vertical loops]
    \label{def:mpdo_vertical_loops}
    \lean{MPOTensor.verticalLoop}
    \lean{MPOTensor.verticalLoopWith}
    \lean{MPOTensor.representativeLoop}
    \lean{MPOTensor.representativeLoop_cast}
    \lean{MPOTensor.verticalLoopWith_apply_eq_trace_mul_verticalTensor}
    \lean{MPOTensor.verticalLoopWith_one}
    \lean{MPOTensor.verticalLoop_ketLeftMul}
    \lean{MPOTensor.verticalLoop_braRightMul}
    \lean{MPOTensor.sectorCompression}
    \lean{MPOTensor.sectorCompression_def}
    \lean{MPOTensor.blockChainTrace}
    \leanok
    \uses{def:mpo_tensor, def:mpo_family, def:vertical_tensor_mpdo,
        def:mpdo_vertical_one_site_actions}
    Closing the vertical indices of a single site of an MPO tensor $M$ gives
    the matrix $E=\sum_i M^{ii}\in\MN{D}$, and inserting a matrix
    $P\in\MN{d}$ into the loop gives
    $E_P=\sum_{i,j}P_{ji}M^{ij}$, the closed loop of the vertically viewed
    product $PM$, equivalently of $MP$ by cyclicity of the single-site
    trace; inserting the identity recovers $E$. On a chain of length $N+1$,
    the compression of the density
    operator by $P$ at the first site is $P_1H^{(N+1)}P_1$ with
    $P_1=P\otimes\Id^{\otimes N}$, and for a matrix
    $E_\alpha\in\MN{D}$ the chain trace with one marked loop is
    $d_\alpha^{(N+1)}=\tr(E_\alpha E^N)$.
    For a representative tensor $M_\alpha$ whose letters are indexed by
    horizontal bond pairs, its closed vertical loop is defined entrywise by
    $(E_\alpha)_{ab}=\tr(M_\alpha^{(a,b)})$.
    In this case the chain trace is the scalar $d_\alpha$ displayed in the
    proof of \cite[Proposition~4.13]{Cirac2016MPDO_arXiv}, resolved by chain
    length: one loop of $M_\alpha$ followed by loops of $M$, closed by the
    horizontal trace.
\end{definition}

\begin{lemma}[The trace of the density operator is a chain of loops]
    \label{lem:mpdo_trace_loop_chain}
    \lean{MPOTensor.sum_evalWord_diag_eq_verticalLoop_pow}
    \lean{MPOTensor.trace_mpo_eq_trace_verticalLoop_pow}
    \leanok
    \uses{def:mpo_family, def:mpdo_vertical_loops}
    Summing the diagonal word evaluations of an MPO tensor over all
    length-$N$ configurations closes the vertical loops site by site,
    \begin{align}
        \sum_{\sigma}
        M^{\sigma_0\sigma_0}\cdots M^{\sigma_{N-1}\sigma_{N-1}}
        &=E^N,
        \notag
    \end{align}
    and hence $\tr H^{(N)}=\tr(E^N)$.
\end{lemma}

\begin{proof}\leanok
    By induction on $N$: splitting off the first site factors the sum as
    $(\sum_a M^{aa})$ times the sum over the remaining $N$
    sites. Taking horizontal traces and summing the diagonal entries of
    $H^{(N)}$ gives the trace identity.
\end{proof}

\begin{theorem}[Trace of a first-site compression]
    \label{thm:mpdo_first_site_compression_trace}
    \lean{MPOTensor.trace_firstSiteMatrix_mul_mpo}
    \lean{MPOTensor.trace_sectorCompression_of_idempotent}
    \leanok
    \uses{def:mpo_family, def:mpdo_vertical_one_site_actions,
        def:mpdo_vertical_loops}
    For every matrix $P\in\MN{d}$ and every chain length $N+1$,
    \begin{align}
        \tr(P_1H^{(N+1)})
        &=\tr(E_P E^N).
        \label{eq:mpdo_first_site_trace}
    \end{align}
    The marked site closes to the loop with $P$ inserted, and the remaining
    $N$ sites close to plain loops. If $P$ is idempotent, then also
    $\tr(P_1H^{(N+1)}P_1)=\tr(E_P E^N)$.
\end{theorem}

\begin{proof}\leanok
    \uses{lem:mpdo_trace_loop_chain, def:mpdo_vertical_one_site_actions}
    For a length-$N$ configuration $\rho$, write
    $E_\rho^{(N)}=M^{\rho_0\rho_0}\cdots M^{\rho_{N-1}\rho_{N-1}}$.
    Expanding at the first site gives
    \begin{align}
        (P_1H^{(N+1)})_{(a,\rho)(a,\rho)}
        &=\sum_i P_{ai} \tr(M^{ia}E_\rho^{(N)}).
        \notag
    \end{align}
    Summing over $a$ and $\rho$ and collecting closes the remaining loops
    by Lemma~\ref{lem:mpdo_trace_loop_chain}. For the
    two-sided compression, first-site actions compose site by site,
    $P_1Q_1=(PQ)_1$, so cyclicity of the trace and idempotency reduce
    $\tr(P_1H^{(N+1)}P_1)$ to $\tr(P_1H^{(N+1)})$.
\end{proof}

\begin{definition}[Sector projectors of the vertical decomposition]
    \label{def:mpdo_sector_projector}
    \lean{MPOTensor.SectorProjectorData}
    \lean{MPOTensor.sectorProjectorData_of_gauge_corner}
    \leanok
    \uses{def:mpdo_vertical_loops, def:vertical_tensor_mpdo}
    A matrix $P\in\MN{d}$ is a sector projector for the weight
    $\mu\in\C$ and the matrix $E_\alpha\in\MN{D}$ if $P$ is an orthogonal
    projector and inserting $P$ into the single-site vertical loop selects
    the sector contribution, $E_P=\mu E_\alpha$.
    In the vertical decomposition
    \begin{align}
        \widetilde M
        &=\bigoplus_{\alpha,k}
          \mu_{\alpha,k}X_{\alpha,k}M_\alpha X_{\alpha,k}^{-1},
        \notag
    \end{align}
    of the vertically viewed tensor, carried by the
    physical space in the proof of
    \cite[Proposition~4.13]{Cirac2016MPDO_arXiv}, the projector
    $P_{\alpha,k}$ selecting the sector $(\alpha,k)$ has this property with
    $\mu=\mu_{\alpha,k}$ and
    $(E_\alpha)_{ab}=\tr(M_\alpha^{(a,b)})$, independently of $k$. Indeed,
    cyclicity of trace cancels the internal similarity
    $X_{\alpha,k}M_\alpha^{(a,b)}X_{\alpha,k}^{-1}$. Applied to the physical
    isometry of a grouped normal sector, the gauge-corner identity constructs
    the projector and proves its loop identity.
\end{definition}

\begin{theorem}[A nonzero vertical corner has a nonzero sector compression]
    \label{thm:mpdo_grouped_sector_compression_nonzero}
    \lean{MPOTensor.IsHorizontalCF.exists_sectorCompression_ne_zero_of_corner}
    \lean{MPOTensor.exists_rangeProjection_corner_ne_zero}
    \leanok
    \uses{def:horizontal_cf_mpdo, def:mpdo_vertical_loops,
        def:vertical_tensor_mpdo, thm:representative_one_sub_mp_zero,
        thm:cpsv_normal_mpv_phase_gauge}
    Let $M$ be in normalized BNT-refined horizontal form and let $P$ be a
    matrix for which
    $P\widetilde M_vP$ is nonzero for some vertical letter $v$. Then the
    first-site compression $P_1H^{(N+1)}P_1$ is nonzero for some $N$.
    Under the same normalized BNT-refined horizontal-form hypothesis, suppose
    in addition
    that $A$ is a normal tensor, $V$ is an isometry, and the compression by
    $V$ is a nonzero scalar multiple of an invertible conjugate of $A$. Then
    the range projector $VV^\dagger$ has a nonzero vertical corner and hence a
    nonzero first-site compression.
    This is the separation step used in
    \cite[Proposition~4.13, line 1898]{Cirac2016MPDO_arXiv}.

    \textbf{Scope restriction (BNT-refined horizontal form):}
    the formal separation invokes the stronger normalized BNT-refined
    predicate. The literal canonical-form version remains part of the missing
    Lemma-L and active-refinement step recorded in
    \path{docs/paper-gaps/cpgsv17_vertical_cf_grouping.tex}.
\end{theorem}

\begin{proof}\leanok
    If every compression vanished, its matrix entries would say that the
    first-site insertion by $P$ has the same matrix product vectors as zero.
    Lemma~L on the normalized BNT-refined horizontal form would then force the
    inserted tensor itself to vanish, contrary to the assumed nonzero corner.
    For the range-projector specialization, normality supplies a nonzero
    letter of $A$. A nonzero scalar and an invertible conjugation preserve
    this letter. If every $VV^\dagger$ corner vanished, compressing by
    $V^\dagger$ and $V$ and using $V^\dagger V=1$ would contradict the exact
    gauge-corner identity.
\end{proof}

\begin{theorem}[The sector trace identity]
    \label{thm:mpdo_sector_trace_identity}
    \lean{MPOTensor.SectorProjectorData.trace_sectorCompression}
    \leanok
    \uses{def:mpdo_sector_projector, thm:mpdo_first_site_compression_trace}
    Let $P$ be a sector projector for the weight $\mu$ and the matrix
    $E_\alpha$. Then for every chain length $N+1$,
    \begin{align}
        \tr(P_1H^{(N+1)}P_1)
        &=\mu\,d_\alpha^{(N+1)}.
        \label{eq:mpdo_sector_trace}
    \end{align}
    Once spectral normalization and grouping identify the reducing
    projectors with sector projectors $P_{\alpha,k}$ satisfying
    Definition~\ref{def:mpdo_sector_projector}, this specializes to the
    trace identification
    $\tr(P_{\alpha,k}H^{(N)}P_{\alpha,k})=\mu_{\alpha,k}d_\alpha$ used in
    the proof of \cite[Proposition~4.13]{Cirac2016MPDO_arXiv}; the source
    states it at the particular length where the compression does not vanish.
\end{theorem}

\begin{proof}\leanok
    \uses{thm:mpdo_first_site_compression_trace}
    By Theorem~\ref{thm:mpdo_first_site_compression_trace} the trace of the
    compression is $\tr(E_P E^N)$, and the sector-projector property
    replaces $E_P$ by $\mu E_\alpha$.
\end{proof}

\begin{theorem}[Positivity of the vertical weights]
    \label{thm:mpdo_sector_weight_pos}
    \lean{MPOTensor.SectorProjectorData.weight_mul_blockChainTrace_pos}
    \lean{MPOTensor.blockChainTrace_pos}
    \lean{MPOTensor.sector_weight_pos}
    \leanok
    \uses{def:mpdo, def:mpdo_sector_projector,
        thm:mpdo_sector_trace_identity,
        thm:mpdo_sector_compression_trace_pos}
    Let $M$ generate an MPDO, let $Q$ and $P$ be sector projectors for the
    weights $\mu_{\alpha,1}$ and $\mu_{\alpha,k}$ and a common matrix
    $E_\alpha$, and suppose $\mu_{\alpha,1}>0$. If the compression of
    $H^{(N+1)}$ by $P$ at the first site is nonzero, then
    \begin{align}
        \mu_{\alpha,k}\,d_\alpha^{(N+1)}&>0,
        \notag\\
        d_\alpha^{(N+1)}&>0,
        \notag\\
        \mu_{\alpha,k}&>0.
        \notag
    \end{align}
    Granted such sector projectors, this yields the conclusion of the trace
    argument in the proof of
    \cite[Proposition~4.13]{Cirac2016MPDO_arXiv}: since w.l.o.g.\
    $\mu_{\alpha,1}>0$, the positive trace forces $d_\alpha>0$ and hence
    $\mu_{\alpha,k}>0$.
\end{theorem}

\begin{proof}\leanok
    \uses{thm:mpdo_sector_trace_identity,
        thm:mpdo_sector_compression_trace_pos}
    By Theorem~\ref{thm:mpdo_sector_compression_trace_pos} the nonzero
    compression has positive trace, which
    Theorem~\ref{thm:mpdo_sector_trace_identity} identifies with
    $\mu_{\alpha,k}d_\alpha^{(N+1)}$. In particular
    $d_\alpha^{(N+1)}\ne 0$, so the trace of the compression by the sector
    projector of $(\alpha,1)$ at the same length equals
    $\mu_{\alpha,1}d_\alpha^{(N+1)}\ne 0$; that compression is therefore
    also nonzero, and its trace is a positive real. Dividing by
    $\mu_{\alpha,1}>0$ gives $d_\alpha^{(N+1)}>0$, and dividing the first
    inequality by $d_\alpha^{(N+1)}$ gives $\mu_{\alpha,k}>0$.
\end{proof}

\begin{theorem}[Gram rigidity of normal tensors]
    \label{thm:normal_tensor_gram_rigidity}
    \lean{MPSTensor.IsNormal.eq_smul_one_of_commute}
    \lean{MPSTensor.IsNormal.conjTranspose_mul_self_eq_smul_one_of_commute}
    \lean{MPSTensor.IsNormal.smul_mem_unitaryGroup_of_commute}
    \lean{Matrix.smul_mem_unitaryGroup_of_conjTranspose_mul_self_eq_smul_one}
    \leanok
    \uses{def:normal}
    Let $M_\alpha$ be a normal tensor. Any matrix commuting with every
    matrix of $M_\alpha$ is a scalar multiple of the identity. In
    particular, if $X\ne 0$ and $X^\dagger X$ commutes with every matrix
    of $M_\alpha$, then $X^\dagger X=\omega\Id$ for a necessarily positive
    constant $\omega$, and $\omega^{-1/2}X$ is unitary. This is the
    normalized form
    $X_{\alpha,k}^\dagger X_{\alpha,k}=\omega_{\alpha,k}\Id$ of the
    Gram-matrix proportionality in the proof of
    \cite[Proposition~4.13]{Cirac2016MPDO_arXiv}, with
    $X_{\alpha,1}=\Id$ and
    $U_{\alpha,k}=\omega_{\alpha,k}^{-1/2}X_{\alpha,k}$. Deriving the
    commutation from self-adjointness of the sector compressions is a
    separate step.
\end{theorem}

\begin{proof}
    \leanok
    Commutation with the matrices of $M_\alpha$ extends to all word
    evaluations, and after blocking these span the full matrix algebra, so
    the commuting matrix is central and hence scalar. The Gram matrix
    $X^\dagger X$ is positive semidefinite and nonzero, so its scalar value
    is a positive real $\omega$; dividing $X$ by $\sqrt{\omega}$ makes the
    Gram matrix the identity.
\end{proof}

\begin{theorem}[Reflected marked-chain identity]
    \label{thm:mpdo_reflected_marked_chain}
    \lean{MPOTensor.horizontalSlice}
    \lean{MPOTensor.reflectedAdjoint}
    \lean{MPOTensor.markedChainCoefficient}
    \lean{MPOTensor.IsMPDO.markedChainCoefficient_eq_reflectedAdjoint}
    \leanok
    \uses{def:mpdo, def:vertical_tensor_mpdo,
        thm:mpo_adjoint_reflection}
    Let $M$ generate positive semidefinite operators on every nonempty chain,
    let $A$ be a vertical tensor, and suppose that a matrix $V$ selects the
    vertical corner $V^\dagger\widetilde M_vV=cA^v$, where $c>0$.
    For vertical indices $r,s$, write
    $(\mc{R}(A)^{rs})_{ab}=(A^{(a,b)})_{rs}$, and set
    $(A^\sharp)^v=(A^{v^{\mathrm{op}}})^\dagger$. Then, for every pair of
    equal-length tail words $\sigma,\tau$,
    \begin{align}
        \tr\!\left(\mc{R}(A)^{rs}
            M^{\sigma_1\tau_1}\cdots M^{\sigma_L\tau_L}\right)
        &=
        \tr\!\left(\mc{R}(A^\sharp)^{rs}
            (M^\dagger)^{\sigma_L\tau_L}\cdots
            (M^\dagger)^{\sigma_1\tau_1}\right).
        \label{eq:mpdo_reflected_chain}
    \end{align}
    Thus the adjoint exchanges the oriented horizontal bond pair and reverses
    the unmarked tail. In particular, this is not a same-tail identity and
    does not identify $A$ with $A^\sharp$ letterwise.
\end{theorem}

\begin{proof}
    \leanok
    Compress the first site of the positive operator by $V$. The resulting
    matrix is Hermitian. Expanding one matrix entry gives the marked chain on
    the left and the complex conjugate of the entry with $r,s$ and the
    ket--bra words exchanged. Taking the conjugate transpose reverses the tail
    product, cyclicity of trace returns the marked factor to the first
    position, and positivity of $c$ permits cancellation without a phase.
\end{proof}

\begin{theorem}[Pairwise Gram conjugation of grouped vertical corners]
    \label{thm:mpdo_pairwise_vertical_gram_conjugation}
    \lean{MPOTensor.gramDressing}
    \lean{MPOTensor.cornerGramCoefficients}
    \lean{MPOTensor.IsMPDO.markedChainCoefficient_gramDressing_eq_reflectedAdjoint}
    \lean{MPOTensor.IsMPDO.markedChainCoefficient_gramDressing_eq_of_two_corners}
    \lean{MPOTensor.IsHorizontalCF.gramDressing_eq_of_two_grouped_corners}
    \leanok
    \uses{def:mpdo, def:horizontal_cf_mpdo,
        thm:mpdo_reflected_marked_chain,
        cor:horizontal_cf_linear_marked_separation}
    Let $M$ be in normalized BNT-refined horizontal form and generate positive
    semidefinite operators on every nonempty chain. Suppose two vertical
    corners are positive scalar multiples of similarities of the same tensor
    $A$:
    \begin{align}
        V_X^\dagger\widetilde M^vV_X
        &=c_X X A^vX^{-1},
        &
        c_X&>0,
        \notag\\
        V_Y^\dagger\widetilde M^vV_Y
        &=c_Y Y A^vY^{-1},
        &
        c_Y&>0.
        \notag
    \end{align}
    Then the two Gram conjugations agree for every vertical letter,
    \begin{align}
        (X^\dagger X)A^v(X^\dagger X)^{-1}
        &=(Y^\dagger Y)A^v(Y^\dagger Y)^{-1}.
        \label{eq:mpdo_pairwise_gram_conj}
    \end{align}
    Normality of $A$ is not required for this equality; it enters only in
    the subsequent proportionality of the Gram matrices. This is the
    pairwise conclusion of Figures~7--8 in
    \cite[Proposition~4.13, lines~1909--1919]{Cirac2016MPDO_arXiv}.

    \textbf{Scope restriction (BNT-refined horizontal form):}
    the marked-tensor separation used here is proved only for normalized
    BNT-refined horizontal form; see
    \path{docs/paper-gaps/cpgsv17_vertical_cf_grouping.tex}.
\end{theorem}

\begin{proof}
    \leanok
    \uses{thm:mpdo_reflected_marked_chain,
        cor:horizontal_cf_linear_marked_separation}
    Conjugate the open indices in the reflected marked-chain identity for
    the first corner by $X^\dagger$. Its left side becomes the marked chain
    of $(X^\dagger X)A^v(X^\dagger X)^{-1}$, while the gauge factors on the
    reflected side cancel and leave the marked chain of $A^\sharp$. The
    same calculation with $Y$ gives the identical reflected chain. Thus the
    two Gram-dressed marks have equal closed chains for every tail word.

    These marks lie in the physical-letter span. Indeed, with
    $L_X=V_XX$ and $R_X=V_X(X^{-1})^\dagger$, their coefficients are
    $f^X_{(r,s),(i,j)}
    =c_X^{-1}\overline{(L_X)_{ir}}(R_X)_{js}$, and similarly for $Y$.
    The corner identities show that these linear combinations are precisely
    the horizontal slices of the two Gram-dressed tensors. Physical-letter
    marked separation therefore identifies every slice, hence every vertical
    letter.
\end{proof}

\begin{theorem}[Gram conjugation for an actual grouped sector]
    \label{thm:mpdo_grouped_sector_gram_conjugation}
    \lean{MPOTensor.IsMPDO.grouped_sector_gram_conj_eq}
    \leanok
    \uses{thm:mpdo_pairwise_vertical_gram_conjugation,
        thm:mpdo_vertical_bnt_grouping_with_isometries}
    Consider the grouped vertical decomposition furnished by
    Theorem~\ref{thm:mpdo_vertical_bnt_grouping_with_isometries}. For a
    phase class $\alpha$ and a copy $q$, transport its representative tensor
    to the bond space of that copy. If $X_{\alpha,q}$ is the corresponding
    gauge, then
    \begin{align}
        (X_{\alpha,q}^\dagger X_{\alpha,q})
        A_{\alpha,q}^v
        (X_{\alpha,q}^\dagger X_{\alpha,q})^{-1}
        &=A_{\alpha,q}^v,
        \label{eq:mpdo_grouped_gram_conj}
    \end{align}
    for every vertical letter $v$. This is Figure~8 for the actual witnesses
    of the grouped decomposition.
\end{theorem}

\begin{proof}
    \leanok
    \uses{thm:mpdo_pairwise_vertical_gram_conjugation,
        thm:mpdo_vertical_bnt_grouping_with_isometries}
    Transport the distinguished corner of the phase class along the equality
    of its bond dimension with that of the copy $q$. Its distinguished gauge
    is the identity. Apply
    Theorem~\ref{thm:mpdo_pairwise_vertical_gram_conjugation} to this corner
    and the corner indexed by $q$. The Gram conjugation belonging to the
    distinguished corner is the identity, which gives the stated equality.
\end{proof}

\begin{theorem}[Gram conjugation from the adjoint dressing]
    \label{thm:gram_conj_of_dressed_adjoint}
    \lean{Matrix.gram_conj_eq_of_dressed_target}
    \lean{Matrix.gram_conj_eq_gram_conj_of_common_dressed_target}
    \lean{Matrix.gram_conj_eq_conjTranspose_of_dressed_adjoint}
    \lean{Matrix.gram_conj_eq_gram_conj_of_dressed_adjoint}
    \lean{Matrix.commute_gram_ratio_of_gram_conj_eq}
    \lean{Matrix.commute_gram_of_dressed_adjoint_of_conjTranspose_eq}
    \leanok
    Let $X$ be invertible and suppose that $B$ and a target $C$ satisfy
    \begin{align}
        XBX^{-1}
        &=(X^{-1})^\dagger C X^\dagger.
        \notag
    \end{align}
    Then conjugation by the Gram matrix carries $B$ to $C$,
    \begin{align}
        (X^\dagger X)B(X^\dagger X)^{-1}
        &=C.
        \label{eq:mpdo_gram_target}
    \end{align}
    Hence any two gauges satisfying this conditional identity with the same
    target have equal Gram conjugations, and the ratio of their Gram matrices
    commutes with $B$. This is a useful algebraic route from an abstract common
    target to the relative identity in
    \cite[Proposition~4.13, lines 1909--1919]{Cirac2016MPDO_arXiv}.
    The reflected marked-chain argument in the source does not identify the
    target separately for either sector.
    The same-letter choice $C=B^\dagger$ and the
    further specialization $B^\dagger=B$ are useful corollaries, but the
    latter is not assumed by the source.
\end{theorem}

\begin{proof}
    \leanok
    Multiply the dressed-target identity by $X^\dagger$ on the left and by
    $(X^{-1})^\dagger$ on the right. The left side becomes
    $(X^\dagger X)B(X^\dagger X)^{-1}$ and the right side collapses to
    $C$. Eliminating the common target for two gauges and multiplying by the
    inverse Gram matrix gives the relative commutation relation.
\end{proof}

\begin{figure}[ht]
    \centering
    $X_\alpha^\dagger X_\alpha=Y_\alpha^\dagger Y_\alpha$
    \begin{center}
        % X_alpha^dagger X_alpha M_alpha X_alpha^{-1} X_alpha^{-1 dagger}
        %   = Y_alpha^dagger Y_alpha M_alpha Y_alpha^{-1} Y_alpha^{-1 dagger},
        % i.e. the two Gram conjugations of M_alpha agree
        % (CPSV16 Fig. 8; the four gauge factors keep Fig. 8's order along
        % the wire, dagger-gauge pair before M_alpha, inverse pair after).
        % Ink-to-index: the wire is the vertical matrix index of the
        % vertical-direction reading, open at both ends; the up and down
        % legs of M_alpha are the paper's two horizontal legs (the letter
        % index of the vertical normal tensor); the capsules are the gauge
        % matrices on the wire and carry no legs.  Both sides carry the
        % same boundary signature
        % (virtual-west=1 | virtual-east=1 | physical-up=1 | physical-down=1).
        \begin{tenkz}[physical=updown, compact]
            \tnX{X_\alpha^{\dagger}} & \tnX{X_\alpha} & \tn{M_\alpha} &
            \tnX{X_\alpha^{-1}} & \tnX{X_\alpha^{-1\dagger}}
        \end{tenkz}
        $ = $
        \begin{tenkz}[physical=updown, compact]
            \tnX{Y_\alpha^{\dagger}} & \tnX{Y_\alpha} & \tn{M_\alpha} &
            \tnX{Y_\alpha^{-1}} & \tnX{Y_\alpha^{-1\dagger}}
        \end{tenkz}
    \end{center}
    \caption{Equality of the Gram conjugations for two copies of the same
    vertical normal tensor:
    $(X_{\alpha,k}^{\dagger}X_{\alpha,k})M_\alpha
    (X_{\alpha,k}^{\dagger}X_{\alpha,k})^{-1}
    =(X_{\alpha,1}^{\dagger}X_{\alpha,1})M_\alpha
    (X_{\alpha,1}^{\dagger}X_{\alpha,1})^{-1}$.
    Theorem~\ref{thm:mpdo_pairwise_vertical_gram_conjugation} gives this
    equality from self-adjointness of the corresponding sector compressions
    and horizontal marked separation. Normality is used only in the
    subsequent proportionality argument. This is
    \cite[Proposition~4.13, lines~1915--1921]{Cirac2016MPDO_arXiv}.}
    \label{fig:mpdo_vertical_gauge_gram_comparison}
\end{figure}

\begin{theorem}[Gram-matrix proportionality from relative Gram conjugations]
    \label{thm:eq3_of_dressed_adjoint}
    \lean{MPSTensor.IsNormal.gram_eq_pos_smul_gram_of_gram_conj_eq}
    \lean{MPSTensor.IsNormal.gram_eq_pos_smul_gram_of_common_dressed_target}
    \lean{MPSTensor.IsNormal.smul_mul_nonsing_inv_mem_unitaryGroup_of_common_dressed_target}
    \lean{Matrix.smul_mul_nonsing_inv_mem_unitaryGroup_of_gram_eq_smul}
    \leanok
    \uses{def:normal, thm:gram_conj_of_dressed_adjoint,
        thm:mpdo_pairwise_vertical_gram_conjugation,
        thm:mpdo_grouped_sector_gram_conjugation,
        thm:normal_tensor_gram_rigidity}
    Let $M_\alpha$ be a normal tensor, and let $X$ and $Y$ be invertible
    gauges whose Gram matrices induce the same conjugation on each letter
    $M_\alpha^v$. Then there is a positive real number $\omega$ such that
    \begin{align}
        X^\dagger X
        &=\omega\,Y^\dagger Y,
        \label{eq:mpdo_gram_proportionality}
    \end{align}
    and $\omega^{-1/2}XY^{-1}$ is unitary.
    Theorem~\ref{thm:mpdo_pairwise_vertical_gram_conjugation} gives the
    Figure~8 comparison in
    \cite[Proposition~4.13, lines 1909--1921]{Cirac2016MPDO_arXiv}. This
    comparison supplies precisely the relative hypothesis for a sector and
    the distinguished sector. In the normalization $Y=\Id$, the resulting
    unitary is $\omega^{-1/2}X$. The common-target declarations attached to
    this entry are conditional algebraic corollaries, not the reflected
    marked-chain statement itself.
\end{theorem}

\begin{proof}
    \leanok
    \uses{thm:gram_conj_of_dressed_adjoint, thm:normal_tensor_gram_rigidity}
    By Theorem~\ref{thm:gram_conj_of_dressed_adjoint}, the ratio
    $(Y^\dagger Y)^{-1}(X^\dagger X)$ commutes with every letter of
    $M_\alpha$. Theorem~\ref{thm:normal_tensor_gram_rigidity} makes this
    ratio scalar. When the matrix dimension is positive, positive
    definiteness of the two Gram matrices makes the scalar a positive real
    number $\omega$. In dimension zero the equality is vacuous and one may
    choose $\omega=1$. Conjugating the relative Gram identity by $Y^{-1}$
    shows that $XY^{-1}$ has Gram matrix $\omega\Id$; division by
    $\sqrt{\omega}$ makes it unitary.
\end{proof}

\begin{theorem}[Normalization of the grouped vertical gauges]
    \label{thm:mpdo_grouped_sector_unitary_normalization}
    \lean{MPOTensor.IsMPDO.grouped_sector_gram_eq_pos_smul_one}
    \lean{MPOTensor.IsMPDO.grouped_sector_exists_unitary_normalization}
    \leanok
    \uses{thm:mpdo_grouped_sector_gram_conjugation,
        thm:eq3_of_dressed_adjoint,
        thm:mpdo_vertical_bnt_grouping_with_isometries}
    In the grouped vertical decomposition of a matrix-product density
    operator in normalized BNT-refined horizontal form, let $X_{\alpha,k}$ be
    the gauge of a copy of the normal representative $M_\alpha$. Fix the
    distinguished gauge to $X_{\alpha,1}=\Id$. Then there is a positive real
    number $\omega_{\alpha,k}$ such that
    \begin{align}
        X_{\alpha,k}^\dagger X_{\alpha,k}
        &=\omega_{\alpha,k}\Id,
        \notag\\
        U_{\alpha,k}
        &=\omega_{\alpha,k}^{-1/2}X_{\alpha,k},
        \notag
    \end{align}
    and $U_{\alpha,k}$ is unitary. This is the normalization in
    \cite[Proposition~4.13, lines~1903--1921]{Cirac2016MPDO_arXiv}.

    \textbf{Scope restriction (BNT-refined horizontal form):}
    this completed normalization inherits the stronger formal horizontal
    predicate and does not close the literal canonical-form refinement gap; see
    \path{docs/paper-gaps/cpgsv17_vertical_cf_grouping.tex}.
\end{theorem}

\begin{proof}
    \leanok
    \uses{thm:mpdo_grouped_sector_gram_conjugation,
        thm:eq3_of_dressed_adjoint}
    Transport the normal representative to the bond space of the chosen
    copy. Theorem~\ref{thm:mpdo_grouped_sector_gram_conjugation} states that
    conjugation by $X_{\alpha,k}^\dagger X_{\alpha,k}$ fixes every letter of
    the transported tensor. Apply
    Theorem~\ref{thm:eq3_of_dressed_adjoint} with the second gauge equal to
    the identity. The resulting Gram identity has a positive real scalar,
    and division by its positive square root gives the stated unitary.
\end{proof}

\begin{theorem}[Normalized physical maps of the grouped vertical sectors]
    \label{thm:mpdo_normalized_grouped_sector_maps}
    \lean{MPOTensor.normalizedGroupedSectorMap}
    \lean{MPOTensor.normalizedGroupedSectorMap_isometry}
    \lean{MPOTensor.normalizedGroupedSectorMap_orthogonal}
    \lean{MPOTensor.normalizedGroupedSectorMap_intertwining}
    \lean{MPOTensor.normalizedGroupedSectorMap_corner}
    \lean{MPOTensor.IsMPDO.exists_normalized_grouped_sector_maps}
    \leanok
    \uses{thm:mpdo_grouped_sector_unitary_normalization,
        thm:mpdo_vertical_bnt_grouping_with_isometries}
    In the grouped vertical decomposition, put
    \begin{align}
        Q_{\alpha,q}
        &=\omega_{\alpha,q}^{-1/2}X_{\alpha,q},
        \notag\\
        W_{\alpha,q}
        &=V_{\alpha,q}Q_{\alpha,q},
        \notag
    \end{align}
    and identify the bond space of each copy with that of its chosen
    representative. Then
    \begin{align}
        W_{\alpha,q}^\dagger W_{\alpha,q}
        &=\Id,
        \notag\\
        W_{\alpha,q}^\dagger W_{\beta,p}
        &=0
        \quad((\alpha,q)\ne(\beta,p)).
        \notag
    \end{align}
    Writing $c_{\alpha,q}=\nu_{\alpha,q}\zeta_{\alpha,q}>0$, the normalized
    maps satisfy
    \begin{align}
        \widetilde M_vW_{\alpha,q}
        &=W_{\alpha,q}(c_{\alpha,q}M_\alpha^v),
        \notag
    \end{align}
    and retain the literal reconstruction
    \begin{align}
        \widetilde M_v
        &=\sum_\alpha\sum_q
          W_{\alpha,q}(c_{\alpha,q}M_\alpha^v)
          W_{\alpha,q}^\dagger.
        \label{eq:mpdo_normalized_reconstruction}
    \end{align}
    These are the normalized sector identities in
    \cite[Proposition~4.13, lines~1895--1921]{Cirac2016MPDO_arXiv}.
\end{theorem}

\begin{proof}
    \leanok
    \uses{thm:mpdo_grouped_sector_unitary_normalization,
        thm:mpdo_vertical_bnt_grouping_with_isometries}
    Each $Q_{\alpha,q}$ is unitary, so right multiplication by it preserves
    the isometry equation and the pairwise orthogonality of the physical
    sector maps. In the intertwining relation, the scalar normalization
    commutes with every letter and the adjacent factors
    $X_{\alpha,q}^{-1}X_{\alpha,q}$ cancel. The Gram identity gives
    $X_{\alpha,q}^\dagger
    =\omega_{\alpha,q}X_{\alpha,q}^{-1}$; hence
    \begin{align}
        Q_{\alpha,q}M_\alpha^vQ_{\alpha,q}^\dagger
        &=X_{\alpha,q}M_\alpha^vX_{\alpha,q}^{-1}.
        \notag
    \end{align}
    Substitution into every summand of the grouped reconstruction proves the
    final identity. Transport along the equality of bond dimensions changes
    only the coordinate labels.
\end{proof}
